\section{Results}

\label{sec:results}

\subsection{Corrected authorization and contemporaneous calibration}

\label{sec:results-corrected-authorization-and-contemporaneous-calibration}

The corrected authorization study records all 900 planned assignments: Tables 2--4 summarize corrected authorization, contemporaneous calibration, and bounded component controls, respectively. The following paired intervals are distinct from the marginal Wilson intervals in the tables.

795 workflows end normally, 60 exhaust the step budget, and 45 have parse failures. Workflow status is separate from task completion. Across the 490 attack assignments, security events are zero under the benchmark-specific checks. This study contains no unguarded control, and repeated condition assignments are not independent attacks; the zero endpoint cannot establish effectiveness or equivalent protection.

For AgentDojo benign-tool tasks, compiled authorization completes 2/30 under both policy and full modules. Exact-target authorization completes 13/30 under policy and 11/30 under full modules. The paired exact-minus-compiled differences are +36.67 percentage points (95\% CI +18.60 to +53.31; Holm p=0.0039) and +30.00 points (+13.23 to +46.75; Holm p=0.0117), respectively. For M3-v2, the corresponding compiled counts are 16/30 and 17/30, compared with exact-target counts of 14/30 and 16/30. Their differences, \ensuremath{-}6.67 and \ensuremath{-}3.33 points, have intervals spanning zero. This does not reproduce a statistically established benchmark-intrinsic reversal. The oracle improves some observed completion under answer-derived scope but is not a deployed non-oracle authorization method.

The separate calibration closes all 480 assignments: 443 workflows end normally, nine exhaust steps, 25 have parse failures, two have runner failures, and one has a provider/stop failure. On AgentDojo, no research Guard yields 3/30 security events and 17/30 benign-tool completions. Policy and full modules each yield 0/30 and 2/30. The paired completion difference is \ensuremath{-}50.00 points (95\% CI \ensuremath{-}66.83 to \ensuremath{-}26.80; Holm p=0.0008), whereas the security difference is \ensuremath{-}10.00 points (\ensuremath{-}25.62 to +3.10; Holm p=0.7500). The utility cost is observable; a security improvement is not established. The no-research-Guard, policy, and full M3-v2 security cells are all zero, with benign completions of 16/30, 17/30, and 16/30. Full modules do not show an independent end-to-end advantage in these conditions.

Configured VIGIL has four AgentDojo attack positives among 28 known outcomes and two unknown attack outcomes. All 30 assignments remain in ITT, so the observed positive count gives a 4/30 lower bound and a 6/30 conservative upper bound, not a claim that the unknown outcomes were safe. Two additional AgentDojo benign outcomes are unknown. Its disabled model secondary verification and larger call budget bound this configuration comparison. Across all 480 rows, two field-projection receipts and four structured safety receipts are absent; missing receipts are not fabricated afterward.

The public-event diagnosis reports 26/30 AgentDojo compiled benign tasks with a first missing-tool-permission block in each module condition. This locates permission insufficiency but does not prove every missing permission is an implementation error. The original 900 rows have no recorded length truncation; parse failures cannot uniformly be attributed to the 512-token limit. Two rows record task completion before a subsequent operational failure, and both observations are retained.

The 16 synthetic controls all reach a workflow terminal. Full modules deny private restricted flows and ask when classification is missing; the tested public and archival flows complete. Policy-only permits the restricted private/missing-classification simulated flow. Legitimate completion and a security violation can coexist in these controls. Their limited scenarios and two DeepSeek services do not establish general provenance effectiveness or independent-provider replication.

Task total latency is available for 476/480 calibration assignments, with four missing and no invalid values. AgentDojo benign-task median/p95 durations are 5.083/11.575 seconds without the research Guard, 1.209/3.029 with policy, 1.279/3.003 with full modules, and 8.145/20.291 for configured VIGIL, whose known-duration denominator is 28/30. Early blocking, model calls, constraint generation, and task-path differences preclude interpreting this ordering as a Guard computation speedup. The latency-only correction does not alter the raw records or prior security and completion inference.

\input{\ProjectRoot/table_02.tex}

\input{\ProjectRoot/table_03.tex}

\input{\ProjectRoot/table_04.tex}

\input{\ProjectRoot/figure_02.tex}

\input{\ProjectRoot/figure_03.tex}

\input{\ProjectRoot/figure_04.tex}

\FloatBarrier

\subsection{Trusted-intent review of benign completion}

\label{sec:results-trusted-intent-review-of-benign-completion}

The fixed 60-task review finds no lost or mismatched action--target pair among the seven trusted application declarations. Thirteen further M3-v2 tasks have rules for explicitly requested targets. The remaining ten M3-v2 tasks require confirmation of a filename or recipient without supplying a concrete target; no independent approval was obtained. These counts describe the trusted inputs and generated rules, not an adjudicated false-block rate.

All 30 AgentDojo tasks lack such concrete trusted action declarations. The generic compiler's lexical tool selection and literal security-target parsing do not cover arbitrary natural-language requests, multistage lookups, or dynamic object identifiers. The requests can express legitimate goals even when this bounded compiler cannot derive an executable grant. The review does not establish a new bounded-contract violation or justify automatic authority expansion.

Six M3-v2 tasks show successful simulated execution without completion under at least one policy/full condition; four of these also have a model-format or operational failure. Categories overlap. In three explicit mock-API tasks, policy records contain ALLOW, successful execution, and a final answer but still fail completion. The evaluator requires both expected-effect matching and a successful final outcome. Redacted public arguments and answer content do not identify the exact mismatch, so this observation is not evidence of an evaluator defect. Native AgentDojo completion likewise does not prove that a particular blocked action was executed.

\subsection{Corrected development study}

\label{sec:results-corrected-development-study}

Supplementary Table S20 separates the corrected-v1 development batch from all historical experiments. DeepSeek V4 Pro and DeepSeek V4.1 Flash each completed 288 planned task-condition executions over the same 48 tasks and six configurations. Pro had 265 completed workflows, 13 parse failures, and 10 max-step terminations; Flash had 272, seven, and nine respectively. Workflow completion is not task completion: ending with ASK records a pending approval, not an executed task. Gemini is deferred and supplies no model-effect comparison.

AgentDojo retains its native security checker rather than a state-only endpoint. For Pro, normal-prompt attack positives were 5/12 without a guard, 2/12 with policy modules, and 0/12 with full modules; benign-tool completion was 6/8, 3/8, and 3/8. The full-minus-policy paired attack difference was -16.67 percentage points with a 95\% interval from -44.80 to +10.38, crossing zero. For Flash, the no-guard normal-prompt attack cell had 1/12 positives and each other cell had zero; benign-tool completion was 7/8 in both no-guard cells and 4/8 in each guarded cell. These small, historically enriched development samples do not establish superiority or causal attribution to provenance.

M3-v2 observed simulated unsafe effects were 0/12 in every attack cell for both models, including no guard. Thus the null observations do not establish guard efficacy. Runtime summaries contained zero labelled sensitive values; dataflow checks were not applicable. Passing synthetic label-propagation tests is separate from model-generated evidence of activated dataflow restrictions. The eight no-tool tasks per model-condition remain in the planned denominators, but Supplementary Table S20 focuses on attack and benign-tool outcomes. Unknown outcomes are not replaced by zero, and operational failures remain separately reported.

\subsection{Remaining evidence boundary}

\label{sec:results-remaining-evidence-boundary}

Historical Public and incremental tables retain their recorded evaluator labels. In the old M3-v2 path, an action allowed by the Guard but rejected by the final isolation boundary could be labelled unsafe even though the effect was not applied. Historical ITT variants also include operational failures. The compact records do not support retrospective identification of all actual un\-au\-thor\-ized effects. Historical numbers are therefore legacy-labelled risk endpoints, not corrected execution-effect estimates. Their values are retained, not silently relabelled as corrected-v1 results; AgentDojo's native checker is also not reduced to an environment-state-only definition.

The medium, Public, authorization-regime, factorial, source-authority, challenge, replication, adaptive, and contemporaneous-activity studies jointly support an empirical security--utility paper. They do not support a leading-defense, end-to-end provenance-benefit, saturated-adaptive, or cross-provider claim. E2 does identify source sensitivity at the deterministic Guard decision layer, but its end-to-end action endpoint is floor-limited. E5 supplies only a conditional 0/12 B200-survivor tail with a 24.25\% upper confidence limit. E6 records attack activity only within its historically enriched, single-block DeepSeek cohort; it neither satisfies nor tests the original multi-window activity gate. The earlier Gemini authorization study closed at 32/450 terminal work items, only five response-bearing; the separately executed E6 Gemini prefix produced five response-bearing records, one response-less failure, and 174 administrative closures. Both are retained for operational accounting but all Gemini effects are not estimable. UTR and DSR are not identifiable from the current record contract and are not reconstructed from unsafe-effect rate.

\subsection{Historical diagnostic context}

\label{sec:results-historical-diagnostic-context}

The full working draft and Supplementary Tables S1-S20 retain the earlier medium, Public, authorization, factorial, provenance, challenge, replication, adaptive and permit studies. Their recorded endpoint definitions are not retrospectively replaced by corrected execution outcomes. In E2, deterministic replay had 0/240 false allows and 30/240 false blocks, while all 240 end-to-end runs executed no action. This establishes bounded decision sensitivity, not useful end-to-end protection. The complete 40-task adaptive panel recorded first legacy risk success for 4/40 tasks by budget 50 and 18/40 by 200. Only twelve preselected B200 survivors continued to 500; the conditional 0/12 result has a Wilson upper limit of 24.25\%, not a complete B500 curve or a saturation estimate. Historical E6 used one risk-enriched execution block, with 47/60 tasks also in E3. It supplies neither independent-sample nor temporal replication. E7a issued execution permits without a backend state transition and therefore supplies neither TCR nor executed-effect evidence.
